\section{Introdução}

A implantodontia constitui uma das especialidades odontológicas que mais se beneficia do paradigma dos gêmeos digitais, em razão de duas características intrínsecas e indissociáveis: a exigência de precisão submilimétrica tridimensional no posicionamento de fixações de titânio e a complexidade biomecânica e biológica do fenômeno de osseointegração. \destaque{Enquanto na ortodontia o movimento dentário pode ser ajustado iterativamente ao longo do tratamento, na implantodontia a interface implante-osso é inexoravelmente determinada no momento da instalação cirúrgica.} Esta interface evolui de forma dinâmica durante o processo de reparo tecidual, demandando modelos preditivos sofisticados que capturem tanto a mecânica inicial (estabilidade primária) quanto a intrincada biologia da cicatrização óssea (estabilidade secundária)~\cite{khoshfekr2025digital}. 

A incorporação de tecnologias digitais ao fluxo de planejamento implantodôntico não é, por si só, uma novidade recente. A utilização da tomografia computadorizada de feixe cônico (CBCT) e a subsequente manipulação de dados em softwares de planejamento virtual (\textit{Computer-Aided Implantology} - CAI) encontram-se estabelecidas clinicamente há mais de duas décadas. Contudo, o que distingue e alicerça a abordagem contemporânea baseada estritamente em \textit{Digital Twins} é a capacidade inaudita de transcender a mera visualização estática~\cite{techsoc2025precision}. Esta nova ontologia digital integra dados estáticos de imagem com simulações biomecânicas dinâmicas, biossensores de monitoramento contínuo (\textit{in vivo}) e algoritmos preditivos impulsionados por Inteligência Artificial (IA), criando uma representação cibernética viva, acoplada e evolutiva do sítio implantar~\cite{duggal2026exploring}. 

Nesse contexto epistemológico, o gêmeo digital atua como uma ponte bidirecional contínua entre os espaços físico e virtual. A convergência entre simulação numérica avançada (como o Método dos Elementos Finitos não lineares), manufatura aditiva de alta fidelidade topológica e monitoramento biológico em tempo real está redefinindo os padrões de previsibilidade e sucesso em reabilitações complexas~\cite{mdpi2025convergence}. 

Este capítulo examina de forma exaustiva a aplicação de gêmeos digitais na implantodontia e na cirurgia oral maxilofacial. A estrutura do texto percorre o planejamento virtual de implantes com abordagens estocásticas, a modelagem biomecânica preditiva do comportamento de interfaces biomateriais, a antecipação computacional das cascatas de osseointegração e o campo disruptivo dos implantes inteligentes (instrumentados com nanossensores). Adicionalmente, explora-se a extensão dessas tecnologias na cirurgia oral avançada, com foco em procedimentos ortognáticos, exodontias guiadas por limiares de fratura biomecânica e microcirurgias de reconstrução mandibular complexa.
